\documentclass[11pt]{article}

\usepackage[margin=1in]{geometry}
\usepackage{amsmath}
\usepackage{amssymb}

\begin{document}
	
	\raggedright
	\setlength{\parindent}{0pt}
	
	\section*{Within-Site Variation in Understory Kelp Density}
	
	Understory kelp abundance is measured along six \(30 \times 2\ \mathrm{m}\)
	transects at each site. Each transect therefore surveys an area of
	\(60\ \mathrm{m}^{2}\). For species \(s\), site \(i\), year \(t\), and
	transect \(j\), transect-level density is defined as
	
	\begin{equation}
		d_{s,i,t,j}
		=
		\frac{N_{s,i,t,j}}{60},
	\end{equation}
	
	where \(N_{s,i,t,j}\) is the number of individuals observed along
	transect \(j\). The mean density across the six transects is
	
	\begin{equation}
		D_{s,i,t}
		=
		\frac{1}{6}
		\sum_{j=1}^{6}
		d_{s,i,t,j}.
	\end{equation}
	
	We use the following metrics to characterize variation in understory
	kelp abundance among transects within a site.
	
	\subsection*{Standard deviation}
	
	The sample standard deviation in transect-level density is
	
	\begin{equation}
		SD_{s,i,t}
		=
		\sqrt{
			\frac{1}{5}
			\sum_{j=1}^{6}
			\left(d_{s,i,t,j}-D_{s,i,t}\right)^2
		}.
	\end{equation}
	
	Standard deviation has the same units as density and describes the
	typical amount by which individual transect densities differ from the
	site-level mean. Larger values indicate greater spatial heterogeneity
	in kelp density within the site.
	
	\subsection*{Variance}
	
	The sample variance in transect-level density is
	
	\begin{equation}
		V_{s,i,t}
		=
		\frac{1}{5}
		\sum_{j=1}^{6}
		\left(d_{s,i,t,j}-D_{s,i,t}\right)^2.
	\end{equation}
	
	Variance is the squared standard deviation and provides a formal
	measure of dispersion among transects. Although its squared units make
	it less intuitive than standard deviation, it is useful in statistical
	analyses and for comparing changes in within-site variability through
	time.
	
	\subsection*{Transect occupancy}
	
	Transect occupancy is defined as the proportion of transects on which
	the species is present:
	
	\begin{equation}
		O_{s,i,t}
		=
		\frac{1}{6}
		\sum_{j=1}^{6}
		I\left(N_{s,i,t,j}>0\right),
	\end{equation}
	
	where \(I(\cdot)\) equals 1 when the condition is true and 0 otherwise.
	Occupancy ranges from 0 to 1 and may be multiplied by 100 to express
	the percentage of occupied transects. This metric indicates how
	broadly a species is distributed across the site, independently of its
	mean density.
	
	\subsection*{Variance-to-mean ratio}
	
	Because all transects survey the same area, the variance-to-mean ratio
	can be calculated from the raw transect counts:
	
	\begin{equation}
		VMR_{s,i,t}
		=
		\frac{
			\displaystyle
			\frac{1}{5}
			\sum_{j=1}^{6}
			\left(N_{s,i,t,j}-\overline{N}_{s,i,t}\right)^2
		}{
			\overline{N}_{s,i,t}
		},
	\end{equation}
	
	where
	
	\begin{equation*}
		\overline{N}_{s,i,t}
		=
		\frac{1}{6}
		\sum_{j=1}^{6}
		N_{s,i,t,j}.
	\end{equation*}
	
	The variance-to-mean ratio, also called an index of dispersion,
	characterizes the spatial distribution of counts among transects.
	Values greater than 1 indicate that counts are more aggregated or
	patchy than expected under a Poisson distribution, whereas values near
	1 are consistent with Poisson-like dispersion. The metric is undefined
	when the mean count is zero and should be interpreted cautiously when
	mean abundance is very low.
	
	Together, these metrics distinguish changes in average kelp density
	from changes in its spatial distribution. Standard deviation and
	variance quantify the magnitude of variation among transects,
	occupancy measures the spatial extent of occurrence within the site,
	and the variance-to-mean ratio describes the relative aggregation of
	individuals among transects.
	
\end{document}